\documentclass[a4paper,oneside,10pt]{article}
\usepackage{wallpaper}
\usepackage{float}
\usepackage[utf8]{inputenc}
\usepackage{circuitikz}

\setlength{\textwidth}{16cm} %
\setlength{\textheight}{23cm} %
\setlength{\oddsidemargin}{0cm} %
\setlength{\evensidemargin}{0cm} %
\setlength{\topmargin}{0cm}

\CenterWallPaper{1.01}{bg.pdf}


\begin{document}

\noindent Relatório x: ASSUNTO. \\
DISCIPLINA (INDIQUE SE É PERÍODO INTEGRAL OU NOTURNO)\\
Prof. Gabriel de Freitas Fernandes\\
202x.x\\
\\
\textbf{EQUIPE:} $\_ \_ \_ \_ \_ \_ \_ \_ \_ \_ \_ \_ \_ \_ \_ \_ \_ \_ \_ \_ \_ \_ \_ \_ \_ \_ \_ \_ \_ \_ \_ \_ \_ \_ \_ \_ \_ \_  \_ \_ \_ \_ \_ \_ \_ \_ \_ \_ \_ \_ \_  \_ \_ \_ \_  \_ \_ \_ \_  \_ \_ \_ \_  \_ \_ \_ \_  \_ \_  \_ \_ \_ \_  \_ \_  \_ \_ \_ \_  \_ \_  \_ \_ \_ \_  \_ \_  \_ \_ \_ \_  \_ \_  \_ \_ \_ \_  \_ \_  \_ \_ \_ \_  \_ \_  \_ \_ \_ \_  \_ \_  \_ \_ \_ \_  \_ \_  \_ \_ \_ $\\ \\
$\_  \_ \_  \_ \_ \_ \_  \_ \_  \_ \_ \_ \_  \_ \_  \_ \_ \_ \_  \_ \_  \_ \_ \_ \_  \_ \_  \_ \_ \_ \_  \_ \_  \_ \_ \_ \_  \_ \_  \_ \_ \_ \_  \_ \_  \_ \_ \_ \_  \_ \_  \_ \_ \_ \_  \_ \_  \_ \_ \_ \_  \_ \_  \_ \_ \_ \_  \_ \_  \_ \_ \_ \_  \_ \_  \_ \_ \_ \_  \_ \_  \_ \_ \_ \_  \_ \_  \_ \_ \_ \_  \_ \_  \_ \_ \_ \_  \_ \_  \_ \_ \_ \_  \_ \_  \_ \_ \_ \_  \_ \_  \_ \_ \_ \_  \_ \_  \_ \_ \_ \_  \_ \_  \_ \_ \_ \_  \_ \_  \_ \_ \_ \_  \_$
\\

\section{Introdução}
Descreva o contexto que a prática se insere. Dê um breve adiantamento de tudo que deverá ser encontrado neste documento, por exemplo, comentar de forma geral o propósito da prática, quais ferramentas foram utilizadas, qual o objetivo almejado e comentar brevemente os resultados obtidos.

\section{Pré-Laboratório (Caso se aplique)}

Exponha todas as etapas de pré-laboratório. Caso hajam múltiplas atividades, separe cada uma da seguinte forma:

\subsection{Atividade Pré-Lab 1(Use o nome correto)}

\subsection{Atividade Pré-Lab 2(Use o nome correto)}

\section{Prática (Caso se aplique)}

\subsection{Montagem Experimental}

\subsection{Código}

\subsection{Resultados e Discussão}

\section{Dificuldades Encontradas (Caso se aplique)}

Esta seção deve ser usada para indicar possíveis melhorias no que foi feito devido à possíveis dificuldades encontradas durante a execução da prática.

\section{Anexos(Caso se aplique)}

Caso o relatório deva ser acompanhado de outros arquivos (que serão também enviados via Moodle) indique aqui uma lista destes. Por ex:

\begin{itemize}
    \item Arquivo .c contendo determinado código.
    \item Documentação de algum componente
    \item etc
\end{itemize}

\end{document}